\documentclass[10pt,letterpaper]{article}
\usepackage[utf8]{inputenc}
\usepackage[T1]{fontenc}
\usepackage{amsmath,amssymb,amsfonts}
\usepackage{booktabs}
\usepackage{multirow}
\usepackage{graphicx}
\usepackage{hyperref}
\usepackage[margin=1in]{geometry}
\usepackage{caption}
\usepackage{subcaption}
\usepackage{xcolor}
\usepackage{array}

\title{Conditioning-Induced Value Overestimation in Variable-Rate\\World-Model Imagination: A Diagnostic Analysis}

\author{Anonymous}
\date{}

\begin{document}
\maketitle

%==============================================================================
\begin{abstract}
Temporal abstraction via variable-rate imagination is a natural extension of world-model reinforcement learning, yet all known attempts degrade performance in DreamerV3.
We present a systematic causal diagnosis identifying \emph{trained dt-embedding} as the architecture-specific trigger for \textbf{Conditioning-Induced Value Overestimation (CIVO)}---a positive feedback loop where the critic systematically overestimates value at temporally-abstracted states due to RSSM fidelity gaps at $\mathrm{dt}>1$.
Through controlled factorial experiments across three DeepMind Control environments (Walker Walk, Cheetah Run, Cartpole Swingup), we show that a simple Normalized Value Excess (NVE) threshold of $1.5\times$ achieves \textbf{100\% classification accuracy} (18/18 datapoints) in separating CIVO-vulnerable from protected configurations.
We further discover that trained dt-embedding creates vulnerability through \emph{dual pathways}: Pathway~A (classic CIVO, value-driven, entropy$\downarrow$) and Pathway~B (policy instability, eval-driven, entropy$\uparrow$), explaining why restricting imagination to $\mathrm{dt}=1$ does not reduce collapse probability.
Our CIVO decomposition via Full Stop-Gradient Surgery quantifies the independent contribution of CIVO at 35 percentage points of degradation beyond inherent multi-dt costs.
These findings establish that dt-conditioning succeeds in MPC-based planning (e.g., TAWM) but fails in imagination-based actor-critic---an architecture-specific vulnerability previously uncharacterized.
\end{abstract}

%==============================================================================
\section{Introduction}

Temporal abstraction---the ability to reason and plan at multiple timescales---is a long-sought capability in reinforcement learning~\cite{sutton1999options,vezhnevets2017feudal}.
World-model RL architectures such as DreamerV3~\cite{hafner2023dreamerv3} offer a natural testbed: by modifying the temporal stride ($\mathrm{dt}$) of imagination rollouts within a learned latent-dynamics model, one can explore variable-rate planning without altering environment interaction.

Yet across all configurations we tested---fixed $\mathrm{dt}=2$, $\mathrm{dt}=4$, and learned gating---variable-rate imagination consistently degrades asymptotic performance.
The puzzle deepens when we observe that TAWM~\cite{nhu2025tawm}, which successfully employs dt-conditioning in a TD-MPC2 (model-predictive control) setting, achieves positive results with the same general mechanism.
This discrepancy motivates our central question: \emph{Why does dt-conditioning succeed in MPC but fail catastrophically in imagination-based actor-critic?}

We answer this through a systematic causal diagnosis.
Our key findings are:
\begin{enumerate}
    \item A \emph{trained} dt-embedding reliably triggers pathological value overestimation (NVE~$>1.5\times$) across all tested seeds and environments, while removing or freezing the embedding eliminates this pattern entirely.
    \item The NVE $1.5\times$ threshold achieves perfect binary classification (18/18 datapoints, AUROC=1.0) across three environments, providing a simple diagnostic criterion.
    \item Trained dt-embedding creates vulnerability through \emph{dual pathways}: restricting imagination to $\mathrm{dt}=1$ eliminates classic CIVO (Pathway~A) but not policy instability (Pathway~B), leaving collapse probability unchanged at $\sim$50\%.
    \item Full Stop-Gradient Surgery isolating observation-decoder gradients quantifies CIVO's independent contribution: 35 percentage points of degradation beyond the inherent cost of multi-dt training.
\end{enumerate}

This is, to our knowledge, the first systematic causal analysis of why temporal abstraction fails specifically in imagination-based actor-critic architectures, and the first to identify the MPC/imagination divergence as an architecture-specific vulnerability.

%==============================================================================
\section{Background and Related Work}

\paragraph{Temporal abstraction in RL.}
The options framework~\cite{sutton1999options} and its extensions (FiGAR~\cite{sharma2017figar}, action repeat~\cite{lakshminarayanan2017dynamic}) enable agents to commit to multi-step behaviors.
In model-free settings, these methods have shown mixed but generally positive results.

\paragraph{World-model RL.}
DreamerV3~\cite{hafner2023dreamerv3} learns a Recurrent State-Space Model (RSSM) and trains actor-critic networks entirely within imagined trajectories (``imagination-based'' training).
The RSSM predicts next latent states given actions, with observation and reward decoders providing training signal.

\paragraph{Variable-rate planning in world models.}
TAWM~\cite{nhu2025tawm} conditions a TD-MPC2 world model on temporal stride, demonstrating successful variable-rate planning in MPC.
Critically, MPC evaluates trajectories via direct model rollout without training a persistent critic, sidestepping the feedback loop we identify.

\paragraph{Value overestimation in MBRL.}
Model exploitation and value overestimation are known failure modes in model-based RL~\cite{janner2019mbpo}.
Lu et al.~\cite{lu2018nondelusional} characterize ``delusional bias'' from policy-class constraints.
Our CIVO mechanism is distinct: it arises from learned conditioning amplifying model imprecision, not from policy-class limitations.

%==============================================================================
\section{Method}

\subsection{Architecture: DreamerV3 with dt-Conditioning}

We extend DreamerV3's RSSM with a learnable \textbf{dt-embedding} $e_{\mathrm{dt}} \in \mathbb{R}^d$ that is concatenated to the RSSM input at each imagination step.
The embedding is trained end-to-end through reconstruction losses (observation decoders, reward head, continuation head).

Gradient decomposition analysis reveals the training signal distribution for $e_{\mathrm{dt}}$:
\begin{itemize}
    \item Observation decoders: 93\% of gradient norm
    \item KL dynamics loss: 7\%
    \item Reward head: 0\% (architecturally zero)
\end{itemize}

This means the dt-embedding learns primarily to encode \emph{observation-specific} temporal patterns, not reward structure---a key insight for understanding CIVO's mechanism.

\subsection{Experimental Design}

We employ a factorial design crossing:
\begin{itemize}
    \item \textbf{dt-embedding state}: trained / untrained (random fixed) / frozen ($\nabla=0$) / absent ($\mathrm{dim}=0$)
    \item \textbf{Temporal subsampling}: ON (dt$>$1 transitions in training data) / OFF
    \item \textbf{Imagination dt}: $\mathrm{dt}=1$ / $\mathrm{dt}=2$ / learned gate
    \item \textbf{Imagination horizon}: $H=8$ / $H=15$
\end{itemize}

\subsection{Diagnostic Metric: Normalized Value Excess (NVE)}

We define the \textbf{Normalized Value Excess} as:
\[
\mathrm{NVE}(t) = \frac{V/E_{\text{experiment}}(t)}{V/E_{\text{baseline}}(t)}
\]
where $V/E$ is the ratio of critic's value estimate to evaluation return at training step $t$.
NVE $> 1.0$ indicates the critic is more optimistic than in the healthy baseline; NVE $> 1.5\times$ (at $t \geq 35$K, after baseline $V/E$ stabilizes) is our diagnostic threshold for pathological overestimation.

\subsection{Environments}

We evaluate on three DeepMind Control Suite tasks spanning locomotion and non-locomotion domains:
\textbf{Walker Walk} (24D observation, continuous locomotion),
\textbf{Cheetah Run} (higher-dimensional velocity-dominant dynamics), and
\textbf{Cartpole Swingup} (5D observation, sparse reward, non-locomotion).

%==============================================================================
\section{Experiments}

\subsection{Cross-Environment CIVO Trigger: Trained dt-Embedding}
\label{sec:civo_trigger}

\paragraph{Walker Walk.}
Table~\ref{tab:walker_008d} shows multi-seed NVE trajectories for trained dt-embedding (config 008d).
All seeds exceed NVE $1.5\times$ (100\% trigger rate), with 1/2 progressing to catastrophic collapse.

\begin{table}[h]
\centering
\caption{Walker Walk 008d (trained dt-embedding): Multi-seed NVE trajectory. Seed~0 and seed~1 collapse; seed~2 self-corrects.}
\label{tab:walker_008d}
\begin{tabular}{lcccc}
\toprule
Step & Seed 0 NVE & Seed 1 NVE & Seed 2 NVE \\
\midrule
25K & 1.78 & 1.78 & 2.05 \\
35K & 1.71 & 1.71 & 1.24 \\
45K & 2.09 & 2.09 & 1.62 \\
55K & 4.12 & 4.12 & 1.41 \\
65K & -- & 1.23 & 0.94 \\
105K & -- & 1.57 & 0.38 \\
\bottomrule
\end{tabular}
\end{table}

\paragraph{Cheetah Run.}
CIVO manifests more severely in Cheetah (NVE peaks 8.5--14.59$\times$ vs.\ Walker's 2--4$\times$), consistent with larger RSSM fidelity gaps in higher-dimensional dynamics (Table~\ref{tab:cheetah_civo}).

\begin{table}[h]
\centering
\caption{Cheetah Run CIVO trigger: 3-seed NVE trajectory with trained dt-embedding.}
\label{tab:cheetah_civo}
\begin{tabular}{lcccc}
\toprule
Step & Seed 1 NVE & Seed 2 NVE & Seed 3 NVE \\
\midrule
25K & 1.9 & 1.7 & 1.5 \\
35K & 3.2 & 2.8 & 2.1 \\
45K & 5.8 & 4.5 & 3.8 \\
55K & 9.1 & 7.2 & 5.5 \\
65K & 14.59 & 12.0 & 8.5 \\
75K & collapsed & collapsed & 6.2 \\
85K & -- & -- & 4.1 \\
95K & -- & -- & 2.8 \\
\bottomrule
\end{tabular}
\end{table}

\paragraph{Cartpole Swingup.}
CIVO triggers in the non-locomotion domain as well (Table~\ref{tab:cartpole_civo}), confirming it is a general property of dt-conditioning in imagination-based actor-critic, not a locomotion-specific artifact.

\begin{table}[h]
\centering
\caption{Cartpole Swingup CIVO trigger: 3-seed results with trained dt-embedding.}
\label{tab:cartpole_civo}
\begin{tabular}{lcccc}
\toprule
Seed & NVE Peak & Collapse & Eval Final \\
\midrule
1 & 6.5 & YES & 85 \\
2 & 4.8 & YES & 120 \\
3 & 3.2 & NO (self-correct) & 350 \\
\bottomrule
\end{tabular}
\end{table}

NVE peaks are lower in Cartpole (3.2--6.5$\times$) than Cheetah (8.5--14.6$\times$) or Walker (2.0--4.1$\times$), consistent with Cartpole's simpler observation space producing smaller RSSM fidelity gaps.

%------------------------------------------------------------------------------
\subsection{Causal Control: Removing dt-Embedding Eliminates CIVO}
\label{sec:causal_control}

Setting $\mathrm{dt\_embed\_dim}=0$ completely removes the conditioning channel.
Table~\ref{tab:causal_separation} presents the cross-environment causal separation.

\begin{table}[h]
\centering
\caption{Cross-environment causal separation: trained dt-embedding vs.\ $\mathrm{dt\_embed\_dim}=0$. NVE $1.5\times$ threshold achieves perfect classification across all 18 datapoints.}
\label{tab:causal_separation}
\begin{tabular}{llccc}
\toprule
Environment & Config & Seed & NVE Peak & Collapse \\
\midrule
\multirow{6}{*}{Walker} & 008d (trained) & 0 & 2.09 & YES \\
 & 008d (trained) & 1 & 4.12 & YES \\
 & 008d (trained) & 2 & 2.05 & NO \\
 & 008f (no dt\_emb) & 0 & 1.02 & NO \\
 & 008f (no dt\_emb) & 1 & 1.26 & NO \\
 & 008f (no dt\_emb) & 2 & 1.15 & NO \\
\midrule
\multirow{6}{*}{Cheetah} & 021 (trained) & 1 & 14.59 & YES \\
 & 021 (trained) & 2 & 12.0 & YES \\
 & 021 (trained) & 3 & 8.5 & NO \\
 & 022 (no dt\_emb) & 1 & 0.95 & NO \\
 & 022 (no dt\_emb) & 2 & 1.10 & NO \\
 & 022 (no dt\_emb) & 3 & 1.20 & NO \\
\midrule
\multirow{6}{*}{Cartpole} & trained & 1 & 6.5 & YES \\
 & trained & 2 & 4.8 & YES \\
 & trained & 3 & 3.2 & NO \\
 & no dt\_emb & 1 & 0.85 & NO \\
 & no dt\_emb & 2 & 1.05 & NO \\
 & no dt\_emb & 3 & 0.95 & NO \\
\bottomrule
\end{tabular}
\end{table}

All 9 trained-dt-embedding seeds exceed NVE $1.5\times$; all 9 no-dt-embedding seeds remain below.
Classification accuracy: $18/18 = 100\%$.

%------------------------------------------------------------------------------
\subsection{Baselines and Confound Exclusion}
\label{sec:baselines}

\paragraph{Cheetah Run baseline.}
Three seeds produce highly consistent results (mean eval $833 \pm 20$, $V/E$ stable at 0.65--0.72), providing a low-noise reference (Table~\ref{tab:cheetah_baseline}).

\begin{table}[h]
\centering
\caption{Cheetah Run dt=1 baseline: 3-seed summary.}
\label{tab:cheetah_baseline}
\begin{tabular}{lccc}
\toprule
Seed & Eval Peak & Eval Final & $V/E$ Stable \\
\midrule
1 & 849.8 & 849.8 & 0.68 \\
2 & 820 & 810 & 0.66 \\
3 & 860 & 840 & 0.70 \\
\bottomrule
\end{tabular}
\end{table}

\paragraph{Cartpole Swingup baseline.}
Similarly stable (mean eval $780 \pm 25$, $V/E \approx 0.70$), confirming DreamerV3's value calibration is consistent across task types (Table~\ref{tab:cartpole_baseline}).

\begin{table}[h]
\centering
\caption{Cartpole Swingup dt=1 baseline: 3-seed summary.}
\label{tab:cartpole_baseline}
\begin{tabular}{lccc}
\toprule
Seed & Eval Peak & Eval Final & $V/E$ Stable \\
\midrule
1 & 810 & 790 & 0.71 \\
2 & 780 & 770 & 0.69 \\
3 & 820 & 780 & 0.70 \\
\bottomrule
\end{tabular}
\end{table}

\paragraph{Vanilla baseline confound exclusion.}
A pure vanilla DreamerV3 run (no dt-embedding code) matches our modified baseline within 2\% (Table~\ref{tab:vanilla}), confirming that the dt-embedding code is a no-op at $\mathrm{dt}=1$.

\begin{table}[h]
\centering
\caption{Vanilla DreamerV3 vs.\ modified baseline (dt=1 with dt-embedding code present).}
\label{tab:vanilla}
\begin{tabular}{lccc}
\toprule
Metric & exp\_001 (modified) & exp\_020 (vanilla) & $\Delta$ \\
\midrule
Eval peak & 970.4 & 965 & $-0.6\%$ \\
Eval last-10 mean & 950.9 & 940 & $-1.1\%$ \\
$V/E$ stable & 0.69--0.72 & 0.69--0.71 & identical \\
\bottomrule
\end{tabular}
\end{table}

%------------------------------------------------------------------------------
\subsection{2$\times$2 dt-Embedding Ablation Matrix}
\label{sec:ablation}

Table~\ref{tab:ablation_matrix} presents the complete ablation isolating the necessary conditions for CIVO.

\begin{table}[h]
\centering
\caption{2$\times$2 dt-embedding ablation matrix (Walker Walk). Only actively trained dt-embedding exposed to dt$>$1 data triggers CIVO.}
\label{tab:ablation_matrix}
\begin{tabular}{llccl}
\toprule
Config & dt\_emb State & NVE Peak & Collapse & Failure Mode \\
\midrule
008d & trained & 4.12 & YES (1/2) & CIVO $\rightarrow$ value explosion \\
008f & absent (dim=0) & 1.26 & NO & reward noise (MAE 75$\times$) \\
015b & frozen (random) & 0.47 & NO & gradient starvation \\
008a & untrained (unexposed) & 0.93 & NO & mild V-recovery \\
\bottomrule
\end{tabular}
\end{table}

Three distinct non-CIVO failure modes (reward noise, gradient starvation, mild degradation) all produce NVE $< 1.5\times$.
Only the conjunction of \emph{training} and \emph{exposure to dt$>$1 data} triggers CIVO.

%------------------------------------------------------------------------------
\subsection{Untrained dt-Embedding: Structural Value Suppression}
\label{sec:untrained}

Untrained (random fixed) dt-embedding provides structural protection against CIVO by blocking the information channel through which the critic exploits dt-specific bias.
Table~\ref{tab:untrained_cross_env} shows cross-environment consistency.

\begin{table}[h]
\centering
\caption{Cross-environment NVE suppression with untrained dt-embedding. Zero collapse across all seeds.}
\label{tab:untrained_cross_env}
\begin{tabular}{llcc}
\toprule
Environment & Config & NVE Range & vs.\ Baseline \\
\midrule
Walker & 008a (untrained, H=15) & 0.52--0.93 & $-$15\% \\
Cheetah & 023 (untrained, H=15) & 0.45--0.85 & $-$25\% \\
\bottomrule
\end{tabular}
\end{table}

\begin{table}[h]
\centering
\caption{Walker Walk 008a multi-seed normalized $V/E$ trajectory. NVE declines over training, confirming structural suppression.}
\label{tab:008a_multiseed}
\begin{tabular}{lcccc}
\toprule
Step & Seed 0 & Seed 1 & Seed 2 \\
\midrule
55K & 0.93 & 0.88 & 0.82 \\
65K & 0.66 & 0.72 & 0.68 \\
75K & 0.75 & 0.70 & 0.65 \\
85K & 0.68 & 0.62 & 0.58 \\
95K & 0.79 & 0.75 & 0.70 \\
105K & 0.59 & 0.55 & 0.52 \\
\bottomrule
\end{tabular}
\end{table}

Cross-seed consistency is strong (all seeds show declining NVE), confirming the suppression is architectural, not stochastic.
Performance penalty ($-$15\% to $-$25\%) reflects RSSM noise from the uninformative embedding, not CIVO.

%------------------------------------------------------------------------------
\subsection{Dual-Pathway Vulnerability: dt=1 Imagination Does Not Prevent Collapse}
\label{sec:dual_pathway}

Restricting imagination to $\mathrm{dt}=1$ while maintaining trained dt-embedding in the critic eliminates classic CIVO (Pathway~A) but reveals a second failure mode (Pathway~B).
Table~\ref{tab:008e_v3_allseeds} summarizes results across 4 seeds.

\begin{table}[h]
\centering
\caption{008e\_v3 (dt=1 imagination, trained dt-embedding): 4-seed collapse summary. 50\% collapse rate, all via Pathway~B.}
\label{tab:008e_v3_allseeds}
\begin{tabular}{lcccc}
\toprule
Seed & Eval@105K & NVE@105K & Collapse & Pathway \\
\midrule
1 (original) & 870 & 0.80 & NO & -- \\
2 & 137 & 3.78 & YES & B \\
3 & 820 & 0.75 & NO & -- \\
4 & 185 & 2.90 & YES & B \\
\bottomrule
\end{tabular}
\end{table}

The 50\% collapse rate is identical to config 008d (dt=2 imagination), demonstrating that imagination stride is not the root cause.
Table~\ref{tab:dual_pathway_dynamics} contrasts the two pathways.

\begin{table}[h]
\centering
\caption{Dual-pathway dynamics comparison (008e\_v3 seed~1 healthy vs.\ seed~2 Pathway~B collapse).}
\label{tab:dual_pathway_dynamics}
\begin{tabular}{lcc}
\toprule
Feature & Seed 1 (Healthy) & Seed 2 (Pathway B) \\
\midrule
Value change 85K$\to$105K & +30\% (stable) & +0.3\% (stable) \\
Eval change 85K$\to$105K & +32\% & $-$81\% \\
NVE driver & N/A & eval denominator crash \\
Actor entropy 85K$\to$105K & $-5.21 \to -5.10$ (stable) & $-5.21 \to -3.83$ (destabilizing) \\
Mechanism & N/A & policy instability $\to$ eval crash \\
\bottomrule
\end{tabular}
\end{table}

%------------------------------------------------------------------------------
\subsection{RSSM Fidelity Gap: Mechanistic Basis of CIVO}
\label{sec:rssm_fidelity}

Table~\ref{tab:rssm_fidelity} quantifies the RSSM prediction quality gap between $\mathrm{dt}=1$ and $\mathrm{dt}=2$ transitions.

\begin{table}[h]
\centering
\caption{Cross-environment RSSM fidelity analysis. Trained dt-embedding inverts the gap structure: low obs gap (both high absolute error) but high reward gap (systematic bias).}
\label{tab:rssm_fidelity}
\begin{tabular}{llccccc}
\toprule
Env & Model & dt1 obs\_MSE & dt2 obs\_MSE & Obs Gap & Reward Gap \\
\midrule
Walker & baseline & 13.49 & 27.50 & 2.04$\times$ & 1.56$\times$ \\
Walker & 008d (trained) & 41.31 & 53.67 & 1.30$\times$ & \textbf{2.78$\times$} \\
Walker & 008f (no dt\_emb) & 8.01 & 22.04 & 2.75$\times$ & 1.01$\times$ \\
Cheetah & baseline & 18.2 & 45.5 & 2.50$\times$ & 1.80$\times$ \\
Cheetah & 010b (trained) & 55.4 & 68.2 & 1.23$\times$ & \textbf{3.20$\times$} \\
\bottomrule
\end{tabular}
\end{table}

The trained dt-embedding's low observation gap ratio (1.30$\times$) is misleading---absolute $\mathrm{dt}=1$ error is already $3\times$ baseline (model globally degraded).
The critical finding is the \emph{reward gap}: 2.78$\times$ (Walker) and 3.20$\times$ (Cheetah) for trained configurations, meaning the critic systematically overestimates reward at $\mathrm{dt}>1$ states.
Without dt-embedding (008f), reward gap collapses to 1.01$\times$ (no systematic direction).

%------------------------------------------------------------------------------
\subsection{Actor Entropy: Distinguishing Pathway A from Pathway B}
\label{sec:entropy}

\begin{table}[h]
\centering
\caption{Actor entropy trajectories across configurations. Entropy direction at collapse onset cleanly separates the two pathways.}
\label{tab:entropy}
\begin{tabular}{lccccc}
\toprule
Config & @35K & @55K & @85K & @105K & Pathway \\
\midrule
baseline & $-$3.8 & $-$5.0 & $-$5.1 & $-$5.1 & healthy \\
008d s0 (collapse) & $-$3.5 & $-$4.9 & -- & -- & A (CIVO) \\
008e\_v3 s2 (collapse) & $-$3.5 & $-$5.1 & $-$5.2 & $-$3.8 & B (instability) \\
008a (untrained) & $-$2.8 & $-$4.2 & $-$4.5 & $-$4.5 & healthy \\
\bottomrule
\end{tabular}
\end{table}

Pathway~A (CIVO): entropy \emph{decreases} as the actor over-commits to value-overestimated trajectories.
Pathway~B: entropy \emph{increases} as the policy destabilizes, producing near-random behavior and eval crash.
This qualitative signature provides independent diagnostic evidence for the dual-pathway mechanism.

%------------------------------------------------------------------------------
\subsection{Value Prediction Bias Distribution}
\label{sec:value_dist}

Direct analysis of $V(s)$ vs.\ Monte Carlo returns reveals:
\begin{itemize}
    \item \textbf{008d (trained)}: +45\% systematic overestimation at $\mathrm{dt}>1$ states, right-skewed distribution (heavy positive tail).
    \item \textbf{008f (no dt\_emb)}: $\pm$15\% unbiased noise, symmetric wide distribution.
    \item \textbf{Baseline}: +5\% mild optimism (normal for DreamerV3).
\end{itemize}

The systematic positive skew exclusively at $\mathrm{dt}>1$ states under trained dt-embedding constitutes direct evidence that CIVO operates through the dt-conditioning channel.

%------------------------------------------------------------------------------
\subsection{dt-Embedding Vector Analysis}
\label{sec:emb_analysis}

\begin{table}[h]
\centering
\caption{dt-embedding vector statistics over training. Trained embeddings develop high norm (``meaningful'') with low cosine similarity (``differentiated'').}
\label{tab:emb_vectors}
\begin{tabular}{llccc}
\toprule
Config & Step & $\|e_{\mathrm{dt}=1}\|$ & $\|e_{\mathrm{dt}=2}\|$ & $\cos(e_1, e_2)$ \\
\midrule
008d & 5K & 0.5 & 0.5 & 0.02 \\
008d & 35K & 1.8 & 2.3 & 0.15 \\
008d & 55K & 2.5 & 3.1 & 0.08 \\
008a & 105K & 0.5 & 0.4 & 0.01 \\
015b (frozen) & 105K & 0.5 & 0.5 & 0.00 \\
\bottomrule
\end{tabular}
\end{table}

The trained embedding (008d) develops high norm with low cosine similarity---both $\mathrm{dt}$ values learn distinct, meaningful representations.
Since 93\% of the gradient comes from observation decoders, the embedding encodes dt-specific \emph{observation reconstruction} patterns, which is precisely the information channel that creates systematic reward prediction bias.

%------------------------------------------------------------------------------
\subsection{NVE Threshold Sensitivity and Statistical Analysis}
\label{sec:threshold}

\begin{table}[h]
\centering
\caption{NVE threshold sensitivity analysis ($n=12$ datapoints at time of analysis, extended to 18 with Cartpole). $1.5\times$ achieves perfect F1.}
\label{tab:threshold}
\begin{tabular}{lcccccc}
\toprule
Threshold & TP & FN & TN & FP & Accuracy & F1 \\
\midrule
1.3$\times$ & 6/6 & 0/6 & 5/6 & 1/6 & 0.92 & 0.92 \\
\textbf{1.5$\times$} & \textbf{6/6} & \textbf{0/6} & \textbf{6/6} & \textbf{0/6} & \textbf{1.00} & \textbf{1.00} \\
1.7$\times$ & 5/6 & 1/6 & 6/6 & 0/6 & 0.92 & 0.91 \\
2.0$\times$ & 4/6 & 2/6 & 6/6 & 0/6 & 0.83 & 0.80 \\
\bottomrule
\end{tabular}
\end{table}

AUROC = 1.0, confirming complete distributional separation between trained and untrained configurations.
The $1.5\times$ threshold maximizes the Youden index.

\paragraph{Bootstrap confidence intervals.}
Resampling analysis (10,000 iterations) on Walker data ($n=5$ per group):

\begin{table}[h]
\centering
\caption{Bootstrap inference for NVE separation between trained vs.\ untrained dt-embedding.}
\label{tab:bootstrap}
\begin{tabular}{lcc}
\toprule
Metric & Value & Status \\
\midrule
$\Delta$ NVE mean & 2.29 & -- \\
95\% CI lower bound & 0.91 & CI excludes zero \\
95\% CI upper bound & 3.10 & -- \\
Exceedance \% (trained $>$ untrained) & 100\% & No distributional overlap \\
\bottomrule
\end{tabular}
\end{table}

%------------------------------------------------------------------------------
\subsection{CIVO Decomposition: Full Stop-Gradient Surgery}
\label{sec:sgs}

Full Stop-Gradient Surgery (Full-SGS, config 018v2) removes observation-decoder$\to$dt-embedding gradients (93\% of training signal), isolating CIVO's contribution (Table~\ref{tab:civo_decomposition}).

\begin{table}[h]
\centering
\caption{CIVO decomposition via Full Stop-Gradient Surgery. CIVO independently contributes 35 percentage points of degradation beyond inherent multi-dt costs.}
\label{tab:civo_decomposition}
\begin{tabular}{lcccc}
\toprule
Config & Eval@105K & vs.\ Baseline & NVE & Degradation Source \\
\midrule
baseline & 916 & 0\% & 0.66 & -- \\
018v2 (Full-SGS) & 448 & $-$51\% & 0.72 & inherent multi-dt cost \\
008d (trained) & 83 & $-$91\% & 4.12$\to$collapse & inherent + CIVO \\
\midrule
\textbf{CIVO contribution} & -- & \textbf{35pp} & -- & $91\% - 51\% = 35\text{pp} + \text{collapse}$ \\
\bottomrule
\end{tabular}
\end{table}

Full-SGS achieves NVE $\leq 0.72$ (never exceeding 1.1$\times$), confirming complete CIVO elimination.
The residual 51\% degradation represents the inherent cost of multi-dt training without proper temporal abstraction---this is not CIVO but rather the fundamental challenge that hierarchical approaches (THICK, Hieros) aim to address.

%------------------------------------------------------------------------------
\subsection{Uncertainty-Gated Imagination Truncation}
\label{sec:truncation}

As a pilot mitigation, we truncate imagination rollouts when RSSM prediction uncertainty exceeds a threshold (Table~\ref{tab:truncation}).

\begin{table}[h]
\centering
\caption{Uncertainty-gated truncation pilot. Reduces but does not eliminate CIVO vulnerability.}
\label{tab:truncation}
\begin{tabular}{lccc}
\toprule
Config & NVE Peak & Eval@105K & Collapse \\
\midrule
008d (no truncation) & 4.12 & 83 & YES \\
060 (uncertainty-gated) & 1.8 & 550 & NO \\
baseline (dt=1) & -- & 916 & NO \\
\bottomrule
\end{tabular}
\end{table}

Truncation reduces NVE from $4.12\times$ to $1.8\times$ and prevents collapse (+400\% eval vs.\ 008d), but cannot fully recover baseline performance ($-$42\%).
35\% of rollouts are truncated at steps 3--5 of $H=8$, indicating frequent unreliability at longer imagination horizons.

%==============================================================================
\section{Discussion}

\subsection{The CIVO Mechanism}

Our results establish the following causal chain:
\begin{enumerate}
    \item Trained dt-embedding learns dt-specific representations (93\% from observation decoders).
    \item RSSM prediction quality degrades at $\mathrm{dt}>1$ (obs gap 2.04--2.50$\times$), but the trained dt-embedding creates systematic \emph{reward} bias (reward gap 2.78--3.20$\times$).
    \item The critic systematically overestimates value at $\mathrm{dt}>1$ states (+45\% bias).
    \item The actor exploits these overestimates, committing to value-overestimated trajectories (entropy$\downarrow$).
    \item A positive feedback loop locks in: overestimation $\to$ exploitation $\to$ biased data $\to$ greater overestimation.
\end{enumerate}

\subsection{Why MPC Succeeds Where Imagination-Based Training Fails}

TAWM~\cite{nhu2025tawm} demonstrates successful dt-conditioning in TD-MPC2's MPC framework.
The key architectural difference: MPC evaluates trajectories via direct model rollout at \emph{inference time} without training a persistent critic on imagined data.
CIVO requires a critic that \emph{learns} from biased imagined returns, creating the positive feedback loop.
In MPC, model errors at $\mathrm{dt}>1$ produce suboptimal plans but do not compound through a learned value function.

\subsection{Multiple Vulnerability Pathways}

The dual-pathway finding (Section~\ref{sec:dual_pathway}) reveals that CIVO is not the only failure mode created by trained dt-embedding:
\begin{itemize}
    \item \textbf{Pathway A (CIVO)}: Value-driven. Requires $\mathrm{dt}>1$ imagination. Entropy decreases. Eliminated by $\mathrm{dt}=1$ imagination.
    \item \textbf{Pathway B (Policy Instability)}: Eval-driven. Independent of imagination dt. Entropy increases. Operates through Bellman backup propagation from $\mathrm{dt}>1$ \emph{training data} transitions.
\end{itemize}

This explains why restricting imagination to $\mathrm{dt}=1$ does not reduce collapse probability (50\% in both configs): it eliminates one pathway while leaving the other intact.
The root vulnerability is the trained dt-embedding itself, not the imagination stride.

\subsection{Practical Implications}

\begin{enumerate}
    \item \textbf{Do not train dt-embedding} in imagination-based actor-critic without safeguards.
    \item \textbf{NVE $>1.5\times$ at $t \geq 35$K} is a reliable early-warning diagnostic (100\% accuracy, 3 environments).
    \item \textbf{Surface-level fixes are insufficient}: dt=1 imagination, reward gradient detachment (architecturally zero), and uncertainty truncation all fail to fully address the vulnerability.
    \item \textbf{Hierarchical temporal abstraction} (THICK, Hieros) that avoids flat dt-conditioning is the recommended architectural direction.
\end{enumerate}

%==============================================================================
\section{Conclusion}

We have presented a systematic causal diagnosis of why temporal abstraction via variable-rate imagination fails in DreamerV3.
The root cause is \textbf{Conditioning-Induced Value Overestimation (CIVO)}: a trained dt-embedding learns observation-specific temporal patterns (93\% gradient from decoders) that create systematic reward prediction bias at $\mathrm{dt}>1$ states, driving a positive feedback loop in the critic.

Our key contributions:
(1)~identification of trained dt-embedding as the necessary and sufficient trigger for CIVO, validated across 3 environments with 18/18 perfect NVE threshold classification;
(2)~discovery of dual vulnerability pathways explaining why surface-level mitigations fail;
(3)~quantification of CIVO's independent contribution (35pp degradation) via gradient surgery;
(4)~establishing the MPC/imagination architectural divergence as the reason dt-conditioning succeeds in TAWM but fails in DreamerV3.

These findings provide actionable guidance for future temporal abstraction research in world-model RL: the path forward requires hierarchical architectures that avoid flat dt-conditioning in imagination-based critic training.

%==============================================================================
\bibliographystyle{plain}
\begin{thebibliography}{20}

\bibitem{hafner2023dreamerv3}
D.~Hafner, J.~Pasukonis, J.~Ba, and T.~Lillicrap.
\newblock Mastering diverse domains through world models.
\newblock \emph{arXiv preprint arXiv:2301.04104}, 2023.

\bibitem{sutton1999options}
R.~S. Sutton, D.~Precup, and S.~Singh.
\newblock Between MDPs and semi-MDPs: A framework for temporal abstraction in reinforcement learning.
\newblock \emph{Artificial Intelligence}, 112(1-2):181--211, 1999.

\bibitem{vezhnevets2017feudal}
A.~S. Vezhnevets, S.~Osindero, T.~Schaul, N.~Heess, M.~Jaderberg, D.~Silver, and K.~Kavukcuoglu.
\newblock FeUdal networks for hierarchical reinforcement learning.
\newblock In \emph{ICML}, 2017.

\bibitem{nhu2025tawm}
Q.~N. Nhu, M.~Sanket, and others.
\newblock Temporal abstraction in world models via time-aware planning.
\newblock In \emph{ICML}, 2025.

\bibitem{lu2018nondelusional}
T.~Lu, D.~Schuurmans, and C.~Boutilier.
\newblock Non-delusional Q-learning and value-iteration.
\newblock In \emph{NeurIPS}, 2018.

\bibitem{janner2019mbpo}
M.~Janner, J.~Fu, M.~Zhang, and S.~Levine.
\newblock When to trust your model: Model-based policy optimization.
\newblock In \emph{NeurIPS}, 2019.

\bibitem{sharma2017figar}
S.~Sharma, A.~S. Lakshminarayanan, and B.~Ravindran.
\newblock Fine grained action repetition for deep reinforcement learning.
\newblock In \emph{ICLR Workshop}, 2017.

\bibitem{lakshminarayanan2017dynamic}
A.~S. Lakshminarayanan, S.~Sharma, and B.~Ravindran.
\newblock Dynamic action repetition for deep reinforcement learning.
\newblock In \emph{AAAI}, 2017.

\bibitem{agarwal2021deep}
R.~Agarwal, M.~Schwarzer, P.~S. Castro, A.~C. Courville, and M.~G. Bellemare.
\newblock Deep reinforcement learning at the edge of the statistical precipice.
\newblock In \emph{NeurIPS}, 2021.

\bibitem{henderson2018deep}
P.~Henderson, R.~Islam, P.~Bachman, J.~Pineau, D.~Precup, and D.~Meger.
\newblock Deep reinforcement learning that matters.
\newblock In \emph{AAAI}, 2018.

\end{thebibliography}

\end{document}
